%\documentclass[12pt,oneside]{amsart}
\documentclass{scrartcl}

%\documentclass{article}

\usepackage{amsthm,amsmath,amssymb}
\usepackage[numbers]{natbib}

%\usepackage[colorlinks]{hyperref}
\usepackage{hyperref}
\hypersetup{
    colorlinks,%
    citecolor=black,%
    filecolor=black,%
    linkcolor=black,%
    urlcolor=black
}
\usepackage{hypernat}
\usepackage[top=2.5cm, bottom=2.5cm, left=2.8cm,
right=2.8cm]{geometry}
\usepackage{graphicx}


%\setlength{\parskip}{0pt}
%\setlength{\parsep}{0pt}
%\setlength{\headsep}{0pt}
%\setlength{\topskip}{0pt}
%\setlength{\topmargin}{0pt}
%\setlength{\topsep}{0pt}
%\setlength{\partopsep}{0pt}

\linespread{1}

%\usepackage{marvosym}

%\textwidth = 6.5 in \textheight = 9.2 in \oddsidemargin = 0.0 in
%\evensidemargin = 0.0 in \topmargin = -0.0 in \headheight = 0.0 in
%\headsep = 0.0in \parskip = 0.0in \parindent = 0.0in
%\def\baselinestretch{0.871}


\newtheorem{theorem}{Theorem}[section]
\newtheorem{proposition}[theorem]{Proposition}
\newtheorem{problem}[theorem]{Problem}
\newtheorem{fact}[theorem]{Fact}
\newtheorem{lemma}[theorem]{Lemma}
\newtheorem{defn}[theorem]{Definition}
\newtheorem{corollary}[theorem]{Corollary}
\newtheorem{remark}{Remark}[section]
\newtheorem{example}[theorem]{Example}

\newcommand{\E}{{\mathbb E}}
\newcommand{\se}{{\mathcal{E}}}
\newcommand{\R}{{\mathbb R}}
\renewcommand{\P}{{\mathbb P}}
\newcommand{\ac}{{\mathcal{A}}}
\newcommand{\A}{{\mathcal{A}}}
\newcommand{\B}{{\mathcal{B}}}
\newcommand{\C}{{\mathcal{C}}}
\newcommand{\M}{{\mathcal{M}}}
\newcommand{\Mf}{{\mathfrak{M}}}
\newcommand{\T}{{\mathcal{T}}}
\newcommand{\Z}{{\mathcal{Z}}}
\newcommand{\K}{{\mathcal{K}}}
\newcommand{\lc}{{\mathcal{L}}}
\newcommand{\Ps}{{\mathcal{P}}}
\newcommand{\G}{{\mathcal{G}}}
\newcommand{\pa}{{\mathring{p}}}
\newcommand{\F}{{\cal F}}
\newcommand{\samp}{{\mathcal{X}}}
\newcommand{\X}{{\mathcal{X}}}
\newcommand{\hi}{{\hat{\Phi}}}
\newcommand{\data}{{\text{data}}}
\newcommand{\relu}{{\text{ReLU}}}

\newcommand{\ridge}{{\hat{\beta}_{\text{ridge}}(\lambda)}}
\newcommand{\lasso}{{\hat{\beta}_{\text{lasso}}(\lambda)}}
\newcommand{\ridgemu}{{\hat{\mu}_t^{\text{ridge}}(\lambda)}}
\newcommand{\lassomu}{{\hat{\mu}_t^{\text{lasso}}(\lambda)}}


\newcounter{rcnt}[section]
\renewcommand{\thercnt}{(\roman{rcnt})}

\def\qt#1{\qquad\text{#1}}

\def\argmin{\mathop{\rm argmin}}
\def\argmax{\mathop{\rm argmax}}


\setlength{\parskip}{1.4 \medskipamount}
%\sloppy
%\linespread{1.3}

\begin{document}

\title{STAT 238 - Bayesian Statistics  \\
  Lecture Thirty Four} 
\subtitle{Spring 2026, UC Berkeley}

\author{Aditya Guntuboyina}


\date{22 April 2026}

\maketitle

\section{Hamiltonian Monte Carlo}

There is a given target probability density $\pi$ and our goal is to
construct a Markov chain which satisfies detailed balance with respect
to $\pi$. Given a current value $x$, HMC constructs the next value $y$
by solving the following second order ODE:
\begin{align}\label{hmc_ode}
  \ddot{x}(t) = \nabla \log \pi(x(t)) ~~ \text{ with initialization
  $x(0) = x$ and $\dot{x}(0) = z \sim N(0, I_d)$}. 
\end{align}
The ODE \eqref{hmc_ode} can be alternatively written in the following
first order form:
\begin{align}\label{hmc_firstorder}
  \begin{pmatrix}
    \dot{x}(t) \\ \dot{v}(t)
  \end{pmatrix} =
  \begin{pmatrix}
    v(t) \\ \nabla \log \pi(x(t)) 
  \end{pmatrix} ~\text{with initialization}~
  \begin{pmatrix}
    x(0) \\ v(0)
  \end{pmatrix} =
  \begin{pmatrix}
    x \\ z
  \end{pmatrix}
\end{align}
which can also be written in terms of the Hamiltonian: 
\begin{align}\label{hamiltonian}
  H(x, v) = H(x_1, \dots, x_d, v_1, \dots, v_d) := -\log \pi(x) +
  \frac{1}{2}\|v\|^2 
\end{align}
as 
\begin{align}\label{hmc_form}
    \begin{pmatrix}
    \dot{x}(t) \\ \dot{v}(t)
  \end{pmatrix} =
  \begin{pmatrix}
    \frac{\partial H}{\partial v}(x(t), v(t)) \\ -\frac{\partial
    H}{\partial x}(x(t), v(t))
  \end{pmatrix} ~\text{with initialization}~
  \begin{pmatrix}
    x(0) \\ v(0)
  \end{pmatrix} =
  \begin{pmatrix}
    x \\ z
  \end{pmatrix}
\end{align}
Most authors (see e.g., \citet{neal2011mcmc}) use the formulation
\eqref{hmc_form} to describe HMC. The advantage of the formulation
\eqref{hmc_ode} is that it makes the connection to MALA very clear
(essentially, MALA proposals are obtained by replacing $\nabla \log
\pi(x(t))$ by $\nabla \log \pi(x(0))$).

On the other hand, the main advantage of the Hamiltonian formulation
is that it also works for choices of the Hamiltonian that are
different from \eqref{hamiltonian}. For example, consider the
alternative choices of the Hamiltonian:
\begin{enumerate}
\item $H(x, v) = -\log \pi(x) + v^T M v/2$ for a positive definite
  matrix $v$. Here the dynamics \eqref{hmc_form} become 
  \begin{align*}
    \begin{pmatrix}
      \dot{x}(t) \\ \dot{v}(t)
    \end{pmatrix} =
    \begin{pmatrix}
      M v(t) \\ -\log \pi(x(t))
    \end{pmatrix}. 
  \end{align*}
 because
  $\frac{\partial H}{\partial x} = -\nabla \log \pi(x)$ and
  $\frac{\partial H}{\partial v} = M v$. 
  This formulation of the Hamiltonian is meaningful when the variables
  $x_1, \dots, x_d$ have different scales. One can also make $M$
  depend on $x$. This is related to Riemannian Manifold HMC (see
  \citet{girolami2011riemann}).
 \item $H(x, v) = -\log \pi(x) + \|v\|_1$. Here the dynamics
   \eqref{hmc_form} become
   \begin{align*}
    \begin{pmatrix}
      \dot{x}(t) \\ \dot{v}(t)
    \end{pmatrix} =
    \begin{pmatrix}
      \text{sign}(v(t)) \\ -\log \pi(x(t))
    \end{pmatrix}
   \end{align*}
 because $\frac{\partial H}{\partial v} = \text{sign}(v)$. Also note
 that $v(t) = (v_1(t), \dots, v_d(t))$ and $\text{sign}(v(t))$ is
 interpreted coordinatewise as $(\text{sign}(v_1(t)), \dots,
 \text{sign}(v_d(t)))$. These are interesting dynamics that cannot be
 written in a second order form as in \eqref{hmc_ode}. 
\end{enumerate}

\section{Hamiltonian Dynamics}
Some understanding of Hamiltonian Dynamics will be useful for
HMC. Hamiltonian Dynamics refers to the ODE \eqref{hmc_form}
\begin{align}
  \label{hmc_dyn}
      \begin{pmatrix}
    \dot{x}(t) \\ \dot{v}(t)
  \end{pmatrix} =
  \begin{pmatrix}
    \frac{\partial H}{\partial v}(x(t), v(t)) \\ -\frac{\partial
    H}{\partial x}(x(t), v(t))
  \end{pmatrix} ~\text{with initialization}~
  \begin{pmatrix}
    x(0) \\ v(0)
  \end{pmatrix} =
  \begin{pmatrix}
    x \\ v
  \end{pmatrix}
\end{align}
The above is the same as \eqref{hmc_form} except that we denote the
initial velocity by $v$ (instead of $v$ in \eqref{hmc_form}). In this
section, we will not use the specific form of the Hamiltonian given by
\eqref{hamiltonian}, but we will take the more general form:
\begin{align}\label{Hgen}
  H(x, v) = -\log \pi(x) + K(v)
\end{align}
where $K(v)$ is a symmetric function of $v$ i.e., $K(-v) = K(v)$. 

We will view the Hamiltonian dynamics (run for a fixed time $\sigma$)
as a mapping from $(x, v)$ (the initial state) to a final state
$(x(\sigma), v(\sigma))$. We will denote this mapping by
$T_{\sigma}(x, v)$:
\begin{align*}
  T_{\sigma}(x, v) = (x(\sigma), v(\sigma)). 
\end{align*}
The space of all pairs of position and velocity $(x, v)$ will be
referred to as the phase space. 

The basic properties of $T_{\sigma}(\cdot)$ that we need for HMC are
summarized in this section. For more details, you can refer to
\citet{neal2011mcmc}. 

\subsection{Property One: Invertibility or Reversibility}

The map $T_{\sigma}(x, v)$ is invertible and its inverse can be
written in a simple manner. Let $S(x, v) = (x, -v)$ be the operator
which flips the velocity. Then
\begin{align}\label{Tinv}
  T_{\sigma}^{-1} = S T_{\sigma} S. 
\end{align}
In other words, to compute $T_{\sigma}^{-1}$ at a point $(y, w)$, we
need to take the following three steps:
\begin{enumerate}
\item First, flip the velocity to obtain $(y, -w)$.
\item Second, run the Hamiltonian dynamics for time $\sigma$ starting
  at $(y, -w)$ to obtain $(x, -v)$ at time $\sigma$.
\item Third, flip velocity again to obtain $(x, v)$. 
\end{enumerate}
The formula \eqref{Tinv} has the following simple consequence:
\begin{align}
  \label{invo}
  T_{\sigma} S = (T_{\sigma} S)^{-1}. 
\end{align}
To see this, simply note
\begin{align*}
  (T_{\sigma} S)^{-1} = S^{-1} T_{\sigma}^{-1} = S^{-1} S T_{\sigma} S
  = T_{\sigma} S. 
\end{align*}
A function $g$ for which $g^{-1} = g$ (which is equivalent to $g(g(x))
= x$) is called an \textit{involution}. The above statement is
therefore equivalent to saying that $T_{\sigma} S$ is an
involution. Involutions have been recently used to unify many MCMC
algorithms (see e.g., \citet{glatt2024sacred}).

\subsection{Property Two: Hamiltonian Conservation}

Hamiltonian dynamics conserve the Hamiltonian in the sense that:
\begin{align}\label{ham_const}
  H(x(t), v(t)) = \text{constant}. 
\end{align}
To see this, just note that
\begin{align*}
  \frac{d}{dt} H(x(t), v(t)) &= \sum_{i=1}^d \left( \frac{\partial H}{\partial
                               x_i} \dot{x}_i(t) + \frac{\partial
                               H}{\partial v_i} \dot{v}_i(t) \right) \\
  &= \sum_{i=1}^d \left(\frac{\partial H}{\partial x_i} \frac{\partial
    H}{\partial v_i} + \frac{\partial H}{\partial v_i} \left(-
    \frac{\partial H}{\partial x_i}\right) \right) = \sum_{i=1}^d \left(\frac{\partial H}{\partial x_i} \frac{\partial
    H}{\partial v_i} - \frac{\partial H}{\partial v_i} 
    \frac{\partial H}{\partial x_i} \right) = 0
\end{align*}

\subsection{Property Three: Volume Preservation}
Volume preservation means that if we take a region $R$ in the phase
space $(x, v)$, then
\begin{align}\label{vol_same}
  \text{vol}(R) = \text{vol}(T_{\sigma}(R))
\end{align}
where $T_{\sigma}(R) = \{T_{\sigma}(x, v): (x, v) \in R\}$. 

The volume preservation property \eqref{vol_same} is equivalent to
\begin{align*}
  \text{vol}(R) &= \text{vol}(T_{\sigma}(R)) = \int I\{(y, w) \in
                  T_{\sigma}(R)\} d(y, w) = \int I\{T_{\sigma}^{-1}
                  (y, w) \in R\} d(y, w). 
\end{align*}
Using the change of variable $(x, v) = T_{\sigma}^{-1}(y, w)$ above,
we get
\begin{align*}
  \int I\{T_{\sigma}^{-1}
                  (y, w) \in R\} d(y, w) = \int I\{(x, v) \in R\}
  |\det J T_{\sigma}(x, v)| d(x, v)
\end{align*}
where $J T_{\sigma}$ denotes the Jacobian of $T_{\sigma}$. Thus volume
preservation can be proved by showing that the Jacobian $J T_{\sigma}$
has determinant equal to 1 for all $(x, v)$:
\begin{align}\label{detjac}
  \det J T_{\sigma}(x, v) = 1 ~~ \text{ for all $(x, v)$}. 
\end{align}
Here is a sketch of the proof of \eqref{detjac}. Assume first that
$\sigma$ is small so that 
$T_{\sigma}$ can be well-approximated by a linear
function:
\begin{align*}
  T_{\sigma}(x, v) \approx
  \begin{pmatrix}
    x \\ v
  \end{pmatrix} + \sigma
  \begin{pmatrix}
    \frac{\partial H}{\partial v}(x, v) \\ -    \frac{\partial H}{\partial x}(x, v)
  \end{pmatrix}
\end{align*}
As a result
\begin{align*}
  J T_{\sigma} \approx I + \sigma
  \begin{pmatrix}
    \frac{\partial^2 H}{\partial x \partial v} &  \frac{\partial^2
                                                 H}{\partial v^2}  \\ -\frac{\partial^2
                                                 H}{\partial x^2} & -
                                                                    \frac{\partial^2
                                                                    H}{\partial
                                                                    x
                                                                    \partial
                                                                    v}
  \end{pmatrix}
\end{align*}
which shows that
\begin{align*}
  \det (J T_{\sigma}) \approx \det   \begin{pmatrix}
1 + \sigma    \frac{\partial^2 H}{\partial x \partial v} &  \sigma\frac{\partial^2
                                                 H}{\partial v^2}  \\ -\sigma\frac{\partial^2
                                                 H}{\partial x^2} & 1 -
                                                                    \sigma
                                                                    \frac{\partial^2
                                                                    H}{\partial
                                                                    x
                                                                    \partial
                                                                    v}
  \end{pmatrix} = 1 + O(\sigma^2). 
\end{align*}
If $\sigma$ is not small, then break $[0, \sigma]$ into $N$ intervals
each of length $\sigma/N$, and then decompose $T_{\sigma}$ as $T_N
\circ T_{N-1} \circ \dots \circ T_1$ where $T_j$ is the mapping
corresponding to Hamiltonian dynamics from $t = (j-1)\sigma/N$ to $t =
j\sigma/N$. For each $j$, we use the previous argument to claim $\det
J T_j = 1 + O(\sigma^2/N^2)$. Taking the product over $j$, we get
\begin{align*}
  \det J T_{\sigma} = \prod_{j=1}^N \left(1 + O(\sigma^2/N^2) \right)
  \rightarrow 1
\end{align*}
as $N \rightarrow \infty$. This completes the heuristic argument for
\eqref{detjac}.

\section{Discretization of \eqref{hmc_firstorder}}

The ODE \eqref{hmc_firstorder} cannot be solved exactly for most
$\pi(\cdot)$. So we have to approximately solve it using a
discretization technique. We fix a step size $\epsilon$ and
approximate the ODE \eqref{hmc_firstorder} at $t = 0, \epsilon, 2
\epsilon, \dots$. More specifically, we will construct:
\begin{align*}
  \begin{pmatrix}
    x(0) \\ v(0)
  \end{pmatrix},
  \begin{pmatrix}
    x(\epsilon) \\ v(\epsilon)
  \end{pmatrix},
  \begin{pmatrix}
    x(2 \epsilon) \\ v(2 \epsilon)
  \end{pmatrix}, \dots
\end{align*}
Below we will describe how to obtain $(x(t+\epsilon), v(t +
\epsilon))$ from $(x(t), v(t))$ (this process will then be iteratively
applied starting at $t = 0$).

The standard approach to discretization of the HMC equation
\eqref{hmc_firstorder} is to use the leapfrog discretization. The
leapfrog discretization uses the following formulae to construct
$x(t+\epsilon)$ and $v(t + \epsilon)$ from $x(t), y(t)$.
\begin{equation}
  \label{leapfrog}
  \begin{split}
&    v\left(t + \frac{\epsilon}{2}\right) = v(t) + \frac{\epsilon}{2}
  \nabla \log \pi(x(t)) \\
& x(t + \epsilon) = x(t) + \epsilon v\left(t + \frac{\epsilon}{2}
  \right) \\
& v(t + \epsilon) = v \left(t + \frac{\epsilon}{2} \right) +
  \frac{\epsilon}{2} \nabla \log \pi(x(t + \epsilon))
  \end{split}
\end{equation}
If we denote this map from $(x(t), y(t))$ to $(x(t + \epsilon), y(t +
\epsilon))$ by $T^{\text{disc}}_{\epsilon}(x, v)$, then it is straightforward to
verify that
\begin{equation}\label{leapfrog_exp}
  \begin{split}
  &T_{\epsilon}^{\text{disc}}(x, v)\\ &= \left(x + \epsilon v +
  \frac{\epsilon^2}{2} \nabla \log \pi(x), v + \frac{\epsilon}{2}
  \nabla \log \pi(x) + \frac{\epsilon}{2}\nabla \log \pi \left(x + \epsilon v +
                                        \frac{\epsilon^2}{2} \nabla
                                        \log \pi(x) \right) \right). 
  \end{split}                                        
\end{equation}
Note that the position update above is reminiscent of MALA.

If we apply $N$ leapfrog steps in succession, then the overall mapping
is given by:
\begin{align*}
  T^{\text{disc}, N}_{\epsilon} = T_{\epsilon}^{\text{disc}} \circ \dots \circ
  T_{\epsilon}^{\text{disc}}. 
\end{align*}
This function $T^{\text{disc}, N}_{\epsilon}$ shares the following two properties of
the continuous Hamiltonian dynamics:
\begin{enumerate}
\item \textbf{Invertibility or Time Reversibility}: One can check that $S
  T_{\epsilon}^{\text{disc}, N} S$ serves as the inverse of
  $T_{\epsilon}^{\text{disc}, N}$. This can be verified by proving
  that
  \begin{align*}
    \left(T_{\epsilon}^{\text{disc}} \right)^{-1} = S
    T_{\epsilon}^{\text{disc}} S. 
  \end{align*}
  Given the explicit formula for $ T_{\epsilon}^{\text{disc}} $ in
  \eqref{leapfrog_exp}, this verification is straightforward.
\item \textbf{Volume Preserving}: $T_{\epsilon}^{\text{disc}, N}$ is
  volume-preserving. This can be verified by showing the the
  determinant of the Jacobian of $T_{\epsilon}^{\text{disc}, N}$.
  equals one. This, in turn, can be verified by proving that the
  determinant of the Jacobian of $T_{\epsilon}^{\text{disc}}$
  equals 1. This can be verified directly by using the formula
  \eqref{leapfrog_exp} or by noting (from \eqref{leapfrog}) that
  $T_{\epsilon}^{\text{disc}}$ is the composition of three mappings:
  \begin{align*}
    T_{\epsilon}^{\text{disc}} = A_{\epsilon/2} \circ B_{\epsilon}
    \circ A_{\epsilon/2}
  \end{align*}
  where
  \begin{align*}
    A_{\epsilon/2}(x, v) = \left(x, v + \frac{\epsilon}{2} \nabla \log
    \pi(x)\right) ~~ \text{ and } ~~ B_{\epsilon}(x, v) = (x +
    \epsilon v, v). 
  \end{align*}
  These mappings are very simple and one can directly verify that they
  are volume preserving (i.e., the determinant of their Jacobians
  equals 1). 
\end{enumerate}
While the discretized Hamiltonian dynamics satisfy invertibility (or
time-reversibility) and volume preservation, they do not satisfy
Hamiltonian conservation (unlike continuous Hamiltonian
dynamics). Because of this, a Metropolis acceptance correction has to
be applied when implementing HMC with leapfrog discretization. This
will be described in the next lecture. 

\section{On stationarity of $\pi$ for the continuous Hamiltonian
  Dynamics}

For the continuous Hamiltonian dynamics given by \eqref{hmc_form}, it
turns that if the initial conditions $x$ and $v$ are distributed
independently according to $x \sim \pi$ and $v \sim N(0, I_d)$, then
$x(\sigma)$ and $v(\sigma)$ are also independently distributed as $x(\sigma)
\sim \pi$ and $v(\sigma) \sim N(0, I_d)$ for every $\sigma$. This can
be proved using the continuity equation that we discussed in the last
lecture. This argument is given below.

\subsection{Continuity equation}

Consider the ODE
\begin{align}\label{gen_ode}
  \dot{X}(t) = V(t, X(t)) ~~\text{for $t \geq 0$}. 
\end{align}
Here $X(t) \in \R^d$ and $V(t, \cdot) : \R^d \rightarrow
\R^d$. Suppose we initialize the ODE at $t = 0$ with a random variable
having density $\rho_b$:
\begin{align*}
  X(0) \sim \rho_b. 
\end{align*}
What them is the density of $X(t)$? Let $\rho(t, x)$ denote the
density of $X(t)$ evaluated at $x$. Then $\rho(t, x)$ satisfies the
following PDE:
\begin{align}\label{continuity}
  \frac{\partial}{\partial t} \rho(t, x) = -\nabla \cdot \left(V(t, x)
  \rho(t, x)
  \right) ~~ \text{ with initialization } \rho(0, x) = \rho_b(x), 
\end{align}
where
\begin{align*}
  \nabla \cdot \left(V(t, x) \rho(t, x) \right) = \text{div}(V(t, x)
  \rho(t, x)) := \sum_{i=1}^d \frac{\partial}{\partial x_i} \left(V_i(t, x)
  \rho(t, x) \right)
\end{align*}
The PDE \eqref{continuity} is known as the continuity equation or
Transport equation, Fokker-Planck equation (it is closely related to
the Fokker-Planck and Kolmogorov Forward equations).

The continuity equation has been popular in the recent literature on
generative modeling (see e.g., \citet[Equation (5.2.8), Theorem 5.2.2,
Section B.1.2]{lai2025principles}).

\subsection{Application to Hamiltonian Monte Carlo}

Observe that \eqref{hmc_form} is a special case of \eqref{gen_ode}
with $x$ in \eqref{gen_ode} corresponding to $(x, v)$ and
\begin{align*}
  V(t, x, v) = \begin{pmatrix}
    \frac{\partial H}{\partial v} \\ -\frac{\partial H}{\partial x}
  \end{pmatrix}
\end{align*}
Suppose now that $x \sim \pi$ so that the initial density of
$(x(0), v(0))$  is
\begin{align*}
\rho_b(x, v) = \pi(x) \phi(v) = (2 \pi)^{-d/2} \exp \left(-H(x, v) \right). 
\end{align*}
We then claim that
\begin{align*}
  \rho(t, x, v) = (2 \pi)^{-d/2} \exp(-H(x, v)) ~~ \text{for every $t \geq 0$}
\end{align*}
satisfies the continuity equation \eqref{continuity}. To see this,
simply note that
\begin{align*}
  \nabla \cdot \left(V \rho \right) &= (2 \pi)^{-d/2} \nabla \cdot
                                      \left(e^{-H}
                                      \begin{pmatrix}
                                            \frac{\partial H}{\partial v} \\ -\frac{\partial H}{\partial x}
                                      \end{pmatrix}
                                      \right)
  \\
  &= (2 \pi)^{-d/2} \sum_{i=1}^d \left[\frac{\partial}{\partial x_i}
    \left(e^{-H} \frac{\partial H}{\partial v_i} \right) - \frac{\partial}{\partial v_i}
    \left(e^{-H} \frac{\partial H}{\partial x_i} \right) \right] \\
&=  (2 \pi)^{-d/2} \sum_{i=1}^d \left[e^{-H} \left(\frac{\partial^2
  H}{\partial x_i \partial v_i} - \frac{\partial H}{\partial x_i}
  \frac{\partial H}{\partial v_i} \right) - e^{-H} \left(\frac{\partial^2
  H}{\partial x_i \partial v_i} - \frac{\partial H}{\partial x_i}
  \frac{\partial H}{\partial v_i} \right)\right]  = 0. 
\end{align*}
This implies that $\rho(t, x, v) = \rho_b(x, v)$ satisfies
\eqref{continuity} (observe that $\frac{\partial}{\partial t} \rho(t, x,
v) =  0$ because $\rho(t, x, v)$ is constant in $t$).

This shows that if $x \sim \pi$, then the solution $x(\sigma)$ to
\eqref{hmc_ode} at any time $\sigma$ is distributed according to
$\pi$. This means that $\pi$ is stationary for the Hamiltonian Markov
Chain. 




\begin{thebibliography}{99}

  \bibitem[Neal(2011)]{neal2011mcmc}
Radford M. Neal.
\newblock MCMC using Hamiltonian dynamics.
\newblock In \emph{Handbook of Markov Chain Monte Carlo}, pages 47--95.
\newblock Chapman and Hall/CRC, 2011.

\bibitem[Girolami and Calderhead(2011)]{girolami2011riemann}
Mark Girolami and Ben Calderhead.
\newblock Riemann manifold Langevin and Hamiltonian Monte Carlo methods.
\newblock \emph{Journal of the Royal Statistical Society: Series B (Statistical Methodology)}, 73(2):123--214, 2011.

  \bibitem[Lai et~al.(2025)Lai, Song, Kim, Mitsufuji, and Ermon]{lai2025principles}
Lai, C.-H., Song, Y., Kim, D., Mitsufuji, Y., and Ermon, S. (2025).
\newblock The principles of diffusion models.
\newblock \emph{arXiv preprint arXiv:2510.21890}.

\bibitem[Glatt-Holtz et~al.(2024)]{glatt2024sacred}
Nathan E. Glatt-Holtz, Andrew J. Holbrook, Justin A. Krometis, Cecilia F. Mondaini, and Ami Sheth.
\newblock Sacred and Profane: from the Involutive Theory of MCMC to Helpful Hamiltonian Hacks.
\newblock \emph{arXiv preprint arXiv:2410.17398}, 2024.

%\bibitem[Roberts and Rosenthal(2004)]{RobertsRosenthal2004}
%Roberts, G. O. and Rosenthal, J. S. (2004).
%General state space Markov chains and MCMC algorithms.
%\textit{Probability Surveys}, \textbf{1}, 20--71.

%\bibitem[Meyn and Tweedie(2009)]{MeynTweedie2009}
%Meyn, S. P. and Tweedie, R. L. (2009).
%\textit{Markov Chains and Stochastic Stability}, 2nd ed.
%Cambridge University Press.


  
%\bibitem[Chib and Greenberg(1995)]{chib1995}
%Chib, S. and Greenberg, E. (1995).
%Understanding the Metropolis--Hastings algorithm.
%\textit{The American Statistician}, \textbf{49}, 327--335.


%\bibitem[MacKay and Peto(1995)]{MacKay1995HierarchicalDirichlet}
%D.~J.~C. MacKay and L.~C. Peto,
%\newblock ``A hierarchical Dirichlet language model,''
%\newblock \emph{Natural Language Engineering}, vol.~1,
%no.~3, pp.~289--308, 1995.

%\bibitem[Rubin(1981)]{Rubin1981BayesianBootstrap}
%D.~B. Rubin,
%\newblock ``The Bayesian bootstrap,''
%\newblock \emph{The Annals of Statistics}, vol.~9,
%no.~1, pp.~130--134, 1981.
  
%\bibitem[Fisher(1934)]{Fisher1934}
%R.~A. Fisher,
%\newblock ``Probability, likelihood and quantity of information in the logic of
%uncertain inference,''
%\newblock \emph{Proceedings of the Royal Society of London. Series~A,
%Containing Papers of a Mathematical and Physical Character}, vol.~146,
%no.~856, pp.~1--8, 1934.

%\bibitem[Efron(2010)]{Efron}
%Efron, B. (2010)
%\textit{Large-scale inference: empirical Bayes methods for estimation,
%testing, and prediction}. Cambridge University Press, 2010.

%\bibitem[Shumway and Stoffer(2017)]{ShumwayStoffer}
%Shumway, R. H. and Stoffer, D. S. (2017)
%\textit{Time Series Analysis and Its Applications: With R
%  Examples}. Springer, 4th edition. 
  
%     \bibitem[Jaynes(2003)]{Jaynes} Jaynes, E. T, (2003)
%  \textit{Probability theory: the logic of science}. Cambridge
%  University Press, 2003.
 
%  \bibitem[Sinai(1992)]{Sinai} Sinai, Y. G, (1992)
%  \textit{Probability theory: an introductory course}. Springer
%  Verlag, 1992.  

%\bibitem[{deGroot(1986)}]{DeGrootBlackwell}DeGroot, M. H. (1986)
%       A conversation with David Blackwell.
%\newblock \emph{Statistical Science}, \textbf{1}, 40--53.

%         \bibitem[Mosteller(1987)]{Mosteller} Mosteller, F, (1987)
%  \textit{Fifty challenging problems in probability with
%    solutions}. Courier Corporation, 1987.

%    \bibitem[MacKay(2003)]{MacKay} MacKay, D. J. C., (2003)
%  \textit{Information theory, inference, and learning
%  algorithms}. Cambridge University Press, 2003.

%\bibitem[{Lindley(2013)}]{Lindley}Lindley, D. V. (2013)
%   \textit{Understanding Uncertainty}. John Wiley and sons, 2013.


\end{thebibliography}







\end{document}

%%% Local Variables:
%%% mode: latex
%%% TeX-master: t
%%% End:
